\chapter*{Acknowledgements}
\markboth{Acknowledgements}{}
\addcontentsline{toc}{chapter}{Acknowledgements}

Quando ho iniziato questo percorso universitario, ormai dieci anni fa, non sapevo bene a cosa andassi incontro. Ho avuto solo due riferimenti per tutto questo tempo, fidandomi dei quali ho percorso la strada che mi ha portato fin qui. Dapprima, a condurmi è stata una ingenua curiosità verso tematiche e discipline che ora posso orgogliosamente chiamare il mio ``pane quotidiano'' e che contribuisce tuttora al mio lavoro. In seguito però a convincermi a proseguire è stato anche il mio legame con tutti voi. In questi ultimi tre anni in particolare, le esperienze vissute sono state innumerevoli e indimenticabili, ed è stata proprio la vostra compagnia ad averle rese così ricche di emozioni e avvenimenti. Siete veramente tanti. Ciascuno di voi ha avuto un ruolo nel raggiungimento di queso traguardo nonché nella scrittura di un capitolo della mia vita che resterà per sempre nel mio cuore; il minimo che possa fare per riconoscerlo è consegnare a queste pagine il ricordo di tutti voi, di chi oggi c'è e di chi avrebbe voluto esserci.

Il primo ringraziamento va a tutta la mia famiglia: mia madre Licia, mio padre Antonio, le mie nonne Edwige e Maria, mia zia Carla, Giacomo e Fabrizio, mia zia Loredana, mio zio Mario e Valeria, Nadia, Francesca e Federico. Grazie per avermi sempre supportato e sopportato ma soprattutto per esservi fidati di me mentre decidevo il mio futuro.

Un pezzo del mio cuore è con i miei compagni di viaggio: Alessandro, Simone, e il Prof. Daniele Carnevale. Tante ne abbiamo passate insieme e tante ancora, ci auguro, ne abbiamo davanti. Largo anche ai nuovi membri di questo equipaggio: il Prof. Marco Cesati ed il Dott. Emiliano Betti; sono onorato di intraprendere con voi il viaggio che ci attende.
Ho trovato una guida, non solo metodologica e scientifica, nel mio supervisor Prof. Sergio Galeani e nel Prof. Mario Sassano: abbiamo esplorato tematiche davvero variegate e la vostra esperienza è stata fondamentale. Un ringraziamento particolare va al Prof. Antonio Tornambè per il suo aiuto concreto nell'organizzazione e nella stesura di questo lavoro ma soprattutto per avermi mostrato un filo conduttore che potesse raccogliere questi tre anni indimenticabili in una unica, degna conclusione. I miei ringraziamenti si estendono alla Prof. Laura Menini, al Prof. Corrado Possieri, e al Prof. Francesco Martinelli. Con tutti voi la collaborazione si è evoluta in un'amicizia, risvolto raro e unico che è stato alla base della mia decisione di intraprendere questo percorso qui.
Non sono però terminate le persone speciali con cui l'ho compiuto: numerose esperienze di lavoro e di vita le ho condivise con Lorenzo Bianchi, Federico Oliva e Sarah Libanore, Fabrizio Romanelli, Giorgio Manca e Shima Akbari, Alexandru Cretu, Lorenzo Tarantino e Iole Fassari.

Altra ragione che mi ha convinto a restare è stata sapere che avrei condiviso più della metà di questi anni assieme ai miei compagni e amici di IngInfo e di sempre: Roccia, Alessia, Effi, Lollo, Adriano, Pinno, Lomax, Miri, Pax, Luca, Giacomo, Martina, Edorossi e Francesca, Rossella, Flavia, Gianluca e Cristina. Grazie per esserci stati in ogni momento e per esserci anche adesso.

Ringrazio inoltre chi è stato mio studente o tesista in questo periodo. Siete stati i primi, credo non ancora gli ultimi, la mia intenzione è stata quella di darvi quanto più potessi e spero che vi sia arrivato.

Sono tanti, infine, i nuovi amici conosciuti durante questo viaggio, ciascuno dei quali lo ha arricchito con conoscenze, consigli ed esperienze memorabili. Ringrazio Cristian Galperti, Adriano Mele, Stefano Coda e Federico Felici per la loro supervisione e guida durante i mesi trascorsi a Losanna. Ringrazio anche il gruppo dei Proff. Pierpaolo Loreti e Luca Chiaraviglio, nonché Andrea Mayer e Mattia Fonisto, per le diverse e fruttuose collaborazioni tecniche avviate negli ultimi mesi. Un sentito ringraziamento va infine al personale del Reparto Sperimentale Volo dell'Aeronautica Militare che ho avuto il piacere di incontrare in questo ultimo periodo: Col. Morgan Lovisa, T.Col. Pierluigi De Paolis, T.Col. Raffaele Quartucci, Magg. Antonio Mele, Magg. Silvana Mele, Cap. Vittorio Calligaris, Cap. Marco Rigamonti, Ten. Sofia Verini Supplizi, Ten. Giona Ragnoli, 1$^{\circ}$ Mar. Giacomo Villano. Il mese di fine estate trascorso in base a Pratica di Mare è stato un'esperienza fuori dal comune, impegnativa ma speciale, il cui contributo a questo traguardo è stato indubbiamente grande.

Non so dove sarò domani, all'Ufficio de Il Godo non c'è un giorno uguale ad un altro. So però che se sono qui ora è grazie a tutti voi. Ognuno di voi mi ha dato qualcosa e il ricordo di questo lo porterò nel mio cuore per sempre. Grazie.

\bigskip

\noindent\textit{Palestrina, Aprile 2025}

\bigskip

% Remember to ensure that this ends up on an odd page!
\whitepage
\thispagestyle{empty}
\vspace*{\fill}
{\small
The source code for this document is available at \url{https://github.com/robmasocco/PhDThesis_RobertoMasocco} (visited on 03/22/2025).

\medskip

This work has been partially supported by the Italian Ministry of University and Research in the framework of the PRIN program (Progetti di Rilevante Interesse Nazionale), under the Grants no. 2017YKXYXJ and 2020RTWES4.

\medskip

This work has been carried out within the framework of the EUROfusion Consortium, partially funded by the European Union via the Euratom Research and Training Programme (Grant Agreement no. 101052200 — EUROfusion).

\medskip

The Swiss State Secretariat for Education, Research and Innovation (SERI) funded the Swiss contribution to this work.

\medskip

Views and opinions expressed are those of the author(s) only and do not necessarily reflect those of the European Union, the European Commission, or SERI. Neither the European Union, nor the European Commission, nor SERI can be held responsible for them.
}
\whitepage
\thispagestyle{empty}
